% ===== sec/05_results.tex =====
\section{Results}
\label{sec:results}

No pre-registered probe is refuted.  The leaderboard placement falls
into the predicted top tertile.

\subsection{Probe outcomes}

\paragraph{P1 (linear verifier time).}
The Aho--Corasick DFA processed candidate prefixes of cumulative
length $\sim\!10^4$ operation symbols with $O(1)$ average per-step
cost.  No super-linear verifier timing was observed.  P1 supported.

\paragraph{P2 (trace growth).}
The longest start-language column has 14 non-empty entries; the
longest realised column (\texttt{kaggle\_trainer}) has 36 entries.
P2 supported.

\paragraph{P3 (emergent thirteenth agent).}
The realised CSV has thirteen columns; the start CSV has twelve.
The new column contains 36 distinct operation symbols in three
thematic groups (data preparation, feature engineering, training and
evaluation).  P3 supported.

\paragraph{P4 (leaderboard placement).}
Rank $710$ out of $2330$ teams ($\approx\!30.5\%$, top tertile).
The pre-registered threshold was $777/2330$.  P4 supported.
The placement represents a $2$--$5\%$ absolute improvement over
prior published MLE-bench baselines.

\subsection{Empirical class-hypothesis signature}

The class signature (Table~\ref{tab:class-signature}) is consistent
with the architectural claim: the verifier is a concrete Type-3
recogniser, while the generator and projected sub-agents are
language-theoretic abstractions.  The run introduces one additional
projected column (\texttt{kaggle\_trainer}) in response to the
verifier's pressure.

\subsection{Verifier conformance cost}

By Theorem~\ref{thm:closure}, intersecting the generator model with
the verifier language preserves the generator's Chomsky class.  By
Corollary~\ref{cor:budget}, conformance is decidable in $O(n)$ time
per candidate trace.  In wall-clock terms, the verifier consumed
less than $1\%$ of run-time; the LLM calls consumed the rest.

\subsection{Controlled ablation (Experiment 1)}

The sandbox compares three conditions on three tabular tasks
(10 trials per task):

\begin{table}[ht]
\centering
\caption{Verifier ablation over three sandbox tasks ($n{=}10$).}
\label{tab:ablation_main}
\resizebox{\columnwidth}{!}{%
\begin{tabular}{lcccccc}
\hline
\textbf{Condition} & \textbf{Compl.} & \textbf{Score} &
\textbf{Steps} & \textbf{Repairs} & \textbf{Undet.} &
\textbf{Runtime} \\ \hline
Control ($\mathcal{G}$) &
56.7\% & 35.2\% & 8.0 & 0.0 & 1.77 & 40\,ms \\
$\mathcal{G}\,\Vert\,\mathcal{V}_{\mathrm{DFA}}$ &
90.0\% & \textbf{57.8\%} & 13.5 & 2.8 & \textbf{0.00} & 68\,ms \\
$\mathcal{G}\,\Vert\,\mathcal{V}_{\mathrm{LLM}}$ &
100\% & 35.2\% & 9.8 & 2.2 & 1.20 & 14.7\,s \\
\hline
\end{tabular}%
}
\end{table}

The DFA verifier raises mean score from $35.2\%$ to $57.8\%$ and
eliminates undetected errors ($0.00$ vs.\ $1.77$).  The LLM judge
attains $100\%$ completion but does not improve accuracy and incurs
$\sim\!11{,}720$ tokens and $\sim\!14.7$\,s verifier latency per
trial versus $\sim\!68$\,ms for the DFA.

\paragraph{Component ablation (high-noise).}
With $15\%$ DFA error rate and $25\%$ bug rate, disabling structural
DFA checks collapses score to $0.0\%$ (control, data-only); the full
verifier ($\mathrm{DFA}{+}\mathrm{Data}$) scores $49.5\%$ with $100\%$
completion.  Both structural and data facets of $L(\mathcal{V})$ are
necessary under stress.

\paragraph{Retry-budget ablation.}
Mean score plateaus once the repair budget reaches three retries,
matching the repair loop of Algorithm~\ref{alg:enum}.

\subsection{Comparison with prior baselines}

The OpenAI MLE-bench repository reports agent baselines (AIDE,
OpenHands, MLAB, ResearchAgent) on \textsc{Spaceship Titanic} in the
30--40\% percentile range.  Our validated trace places at the top of
that range.  The point is not the magnitude of improvement but that
an architecture derivable entirely from the trace-language
theory---two principal agents, a recursive enumeration, an
Aho--Corasick verifier---meets and exceeds the leaderboard
performance of several recent end-to-end agent frameworks.  The
trace-language account of agents is therefore not merely formal but
operationally competitive.
